\documentclass[aspectratio=169]{beamer}

% ===============================
% Theme and Color Settings - MINIMALIST UNIQUE DESIGN
% ===============================
\usepackage{tikz}
\usetikzlibrary{shapes.geometric, arrows.meta, positioning, fit, backgrounds, calc}

% Minimalist theme base
\usetheme{default}
\setbeamertemplate{navigation symbols}{}

% Simple, refined color palette
\definecolor{primary}{RGB}{0, 102, 102}        % Deep teal
\definecolor{accent}{RGB}{230, 126, 34}        % Warm orange accent
\definecolor{textcolor}{RGB}{40, 40, 40}       % Soft black
\definecolor{lightgray}{RGB}{245, 245, 245}    % Very light gray
\definecolor{mediumgray}{RGB}{200, 200, 200}   % Medium gray for subtle lines

% Background - pure white for minimalism
\setbeamercolor{background canvas}{bg=white}
\setbeamercolor{normal text}{fg=textcolor}

% Minimalist frame title - simple line accent
\setbeamercolor{frametitle}{fg=primary,bg=white}
\setbeamerfont{frametitle}{size=\LARGE,series=\bfseries}

\setbeamertemplate{frametitle}{
  \vspace{0.4cm}
  \begin{beamercolorbox}[wd=\paperwidth,leftskip=0.5cm,rightskip=0.5cm]{frametitle}
    \usebeamerfont{frametitle}\insertframetitle
    \vspace{0.15cm}
  \end{beamercolorbox}
  \nointerlineskip
  \vspace{-0.05cm}
  \hspace{0.5cm}{\color{accent}\rule{3.5cm}{0.12cm}}
  \vspace{0.25cm}
}

% Text colors
\setbeamercolor{structure}{fg=primary}
\setbeamercolor{alerted text}{fg=accent}

% Minimalist blocks
\setbeamercolor{block title}{bg=primary,fg=white}
\setbeamercolor{block body}{bg=lightgray,fg=textcolor}
\setbeamercolor{block title alerted}{bg=accent,fg=white}
\setbeamercolor{block body alerted}{bg=accent!10,fg=textcolor}

\setbeamertemplate{blocks}[default]

% Simple bullet points
\setbeamertemplate{itemize item}{\small\raise0.5pt\hbox{$\bullet$}}
\setbeamertemplate{itemize subitem}{\small\raise0.5pt\hbox{$\circ$}}
\setbeamertemplate{itemize subsubitem}{\small\raise0.5pt\hbox{--}}

% ===============================
% Additional Packages
% ===============================
\usepackage[utf8]{inputenc}
\usepackage[T1]{fontenc}
\usepackage[french]{babel}
\usepackage{graphicx}
\usepackage{amsmath}
\usepackage{booktabs}

% ===============================
% Presentation Information
% ===============================
\title[Détection d'Anomalies MAS]{Détection d'Anomalies par Système Multi-Agents Distribué}
\subtitle{Edge AI pour Réseaux IoT}
\author[ESSEBYITY A. \& OUARIR H.]{Abderrahman ESSEBYITY \and Hamza OUARIR}
\institute[Master SDA]{Master Science des Données et Analytique}
\date{2025 -- 2026}

% Minimalist footer
\setbeamertemplate{footline}{
  \vspace{0.1cm}
  {\color{mediumgray}\hrule height 0.5pt}
  \vspace{0.1cm}
  \leavevmode%
  \hbox{%
  \begin{beamercolorbox}[wd=.5\paperwidth,ht=2.5ex,dp=1ex,left,leftskip=0.5cm]{author in head/foot}%
    \usebeamerfont{author in head/foot}\color{textcolor}\insertshortauthor
  \end{beamercolorbox}%
  \begin{beamercolorbox}[wd=.5\paperwidth,ht=2.5ex,dp=1ex,right,rightskip=0.5cm]{page in head/foot}%
    \usebeamerfont{page in head/foot}\color{textcolor}
    \insertframenumber{} / \inserttotalframenumber
  \end{beamercolorbox}}%
  \vskip0pt%
}

\begin{document}

% ===============================
% Title Slide
% ===============================
{
\setbeamertemplate{footline}{}
\begin{frame}[plain]
    \vspace{0.2cm}
    \begin{center}
        {\color{accent}\rule{0.7\paperwidth}{0.15cm}}
    \end{center}
    \vspace{1.2cm}
    
    \begin{center}
        {\Huge\bfseries\color{primary} Détection d'Anomalies}\\[0.25cm]
        {\Huge\bfseries\color{primary} par Système Multi-Agents}\\[0.7cm]
        {\color{accent}\rule{0.35\textwidth}{0.12cm}}\\[0.5cm]
        {\large\color{textcolor} Edge AI et Apprentissage en Ligne\\pour Réseaux de Capteurs IoT}
    \end{center}
    
    \vspace{1cm}
    
    \begin{center}
        {\large\textbf{Abderrahman ESSEBYITY} \quad \textbf{\&} \quad \textbf{Hamza OUARIR}}\\[0.35cm]
        {\normalsize Encadré par: \textbf{Saif Islam Bouderba}}\\[0.5cm]
        {\normalsize Master Science des Données et Analytique}\\
        {\normalsize Intelligence Artificielle Distribuée}\\[0.3cm]
        {\small \textbf{2025 -- 2026}}
    \end{center}
    
    \vspace{0.5cm}
    \begin{center}
        {\color{accent}\rule{0.7\paperwidth}{0.15cm}}
    \end{center}
\end{frame}
}

% ===============================
% SLIDE 1: Le Problème
% ===============================
\begin{frame}{Le Défi: Surveillance de Milliers de Capteurs IoT}
    \begin{columns}[T]
        \begin{column}{0.48\textwidth}
            \textbf{Approche Centralisée:}
            \begin{itemize}
                \item[$-$] Surcharge réseau massive
                \item[$-$] Latence élevée (RTT)
                \item[$-$] Point de défaillance unique
                \item[$-$] Coûts prohibitifs
            \end{itemize}
        \end{column}
        
        \begin{column}{0.48\textwidth}
            \textbf{Notre Solution:}
            \begin{itemize}
                \item[$+$] Intelligence en périphérie (Edge AI)
                \item[$+$] Détection temps réel
                \item[$+$] Validation collaborative
                \item[$+$] Scalabilité horizontale
            \end{itemize}
        \end{column}
    \end{columns}
    
    \vspace{0.5cm}
    
    \begin{alertblock}{Question de Recherche}
        Comment détecter des anomalies de manière \textbf{distribuée}, \textbf{autonome} et \textbf{collaborative} dans un réseau de capteurs IoT ?
    \end{alertblock}
\end{frame}

% ===============================
% SLIDE 2: Notre Solution
% ===============================
\begin{frame}{Notre Solution: L'IA Distribuée Collaborative}
    \begin{block}{Principe Clé}
        Déployer des \textbf{Agents Intelligents} directement sur chaque capteur
    \end{block}
    
    \vspace{0.5cm}
    
    \textbf{Les 3 Étapes du Système:}
    
    \begin{enumerate}
        \item \textbf{Détection Locale (O(1))}
        \begin{itemize}
            \item Apprentissage en ligne \textbf{River Half-Space Trees}
            \item Complexité mémoire constante, pas d'historique
        \end{itemize}
        
        \vspace{0.2cm}
        
        \item \textbf{Consensus Distribué}
        \begin{itemize}
            \item Vote majoritaire entre voisins
            \item Réduction des faux positifs de \textbf{37\%}
        \end{itemize}
        
        \vspace{0.2cm}
        
        \item \textbf{Système de Confiance}
        \begin{itemize}
            \item Score de confiance dynamique par agent
            \item Pondération des votes selon l'historique
        \end{itemize}
    \end{enumerate}
\end{frame}

% ===============================
% SLIDE 3a: Topologie Réseau
% ===============================
\begin{frame}{Architecture du Système: Topologie}
    \begin{figure}
        \centering
        \scalebox{0.6}{
        \begin{tikzpicture}[
            agent/.style={circle, draw=teal!70, line width=1.2pt, fill=teal!15, minimum size=1.1cm, font=\normalsize\sffamily\bfseries},
            coordinator/.style={circle, draw=orange!70, line width=1.5pt, fill=orange!20, minimum size=1.4cm, font=\large\sffamily\bfseries},
            comm/.style={->, line width=1pt, color=teal!60},
            alert/.style={->, dashed, line width=1.2pt, color=red!60}
        ]
        
        % Coordinator in center
        \node[coordinator] (coord) at (0,0) {COORD};
        \node[below=0.08cm of coord, font=\footnotesize, color=orange!70] {Observateur};
        
        % Agents in RING topology
        \foreach \i in {1,...,8} {
            \pgfmathsetmacro\angle{360/8*\i + 90}
            \node[agent] (a\i) at ({2.6*cos(\angle)},{2.6*sin(\angle)}) {A\i};
        }
        
        % Ring connections
        \foreach \i in {1,...,7} {
            \pgfmathtruncatemacro\next{\i+1}
            \draw[comm] (a\i) -- (a\next);
        }
        \draw[comm] (a8) -- (a1);
        
        % Alert connections - all agents to coordinator
        \foreach \i in {1,...,8} {
            \draw[alert] (a\i) -- (coord);
        }
        
        \end{tikzpicture}
        }
    \end{figure}
    
    \begin{columns}[t]
        \begin{column}{0.48\textwidth}
            \begin{block}{Topologie Anneau}
                {\small
                \begin{itemize}
                    \item 20 agents distribués
                    \item 2 voisins par agent
                    \item Communication locale
                \end{itemize}
                }
            \end{block}
        \end{column}
        
        \begin{column}{0.48\textwidth}
            \begin{block}{Coordinateur}
                {\small
                \begin{itemize}
                    \item Observateur passif
                    \item Collecte les alertes
                    \item Pas de contrôle central
                \end{itemize}
                }
            \end{block}
        \end{column}
    \end{columns}
\end{frame}

% ===============================
% SLIDE 3b: Pipeline de Traitement
% ===============================
\begin{frame}{Architecture du Système: Pipeline}
    \begin{figure}
        \centering
        \scalebox{0.8}{
        \begin{tikzpicture}[
            component/.style={rectangle, draw=black!60, line width=1.5pt, text width=3cm, text centered, rounded corners=5pt, minimum height=1.3cm, font=\normalsize\sffamily\bfseries},
            arrow/.style={-Stealth, line width=2pt, color=teal!70}
        ]
        
        % Components
        \node[component, fill=blue!15, draw=blue!60] (sensor) at (0,0) {1. Capteur IoT};
        \node[component, fill=green!15, draw=green!60] (ml) at (4.5,0) {2. ML en Ligne\\{\small River HST}};
        \node[component, fill=purple!15, draw=purple!60] (consensus) at (9,0) {3. Consensus\\{\small Vote}};
        \node[component, fill=red!15, draw=red!60] (trust) at (0,-3) {4. Confiance\\{\small Score}};
        
        % Arrows
        \draw[arrow] (sensor) -- (ml);
        \draw[arrow] (ml) -- (consensus);
        \draw[arrow] (consensus) |- (trust);
        
        \end{tikzpicture}
        }
    \end{figure}
    
    \vspace{0.4cm}
    
    \begin{block}{Flux de Traitement Distribué}
        Chaque agent exécute le pipeline complet de manière autonome: acquisition → détection → consensus → confiance
    \end{block}
\end{frame}

% ===============================
% SLIDE 4: Technologies
% ===============================
\begin{frame}{Technologies Clés}
    \begin{columns}[T]
        \begin{column}{0.48\textwidth}
            \textbf{Apprentissage en Ligne:}
            \begin{itemize}
                \item \textbf{River Half-Space Trees}
                \item Complexité O(1) en mémoire
                \item Adaptation au concept drift
                \item Pas de réentraînement nécessaire
            \end{itemize}
            
            \vspace{0.4cm}
            
            \begin{block}{Avantage}
                Contrairement au batch ML, notre système a une empreinte mémoire \textbf{constante}
            \end{block}
        \end{column}
        
        \begin{column}{0.48\textwidth}
            \textbf{Communication Distribuée:}
            \begin{itemize}
                \item \textbf{SPADE Framework} (XMPP)
                \item Messages asynchrones
                \item Latence < 200ms
            \end{itemize}
            
            \vspace{0.4cm}
            
            \textbf{Infrastructure:}
            \begin{itemize}
                \item Docker Compose
                \item Prosody XMPP Server
                \item 20 agents concurrents
            \end{itemize}
        \end{column}
    \end{columns}
\end{frame}

% ===============================
% DÉMO SECTION (3 slides)
% ===============================
\begin{frame}{Dashboard ``War Room'' - Interface de Contrôle}
    \begin{center}
        {\LARGE\textbf{Démonstration Système}}
    \end{center}
    
    \vspace{0.5cm}
    
    \begin{block}{Fonctionnalités du Dashboard}
        \begin{itemize}
            \item \textbf{Mission Control}: Démarrage/Arrêt simulations
            \item \textbf{Topologie Live}: Graphe dynamique des connexions
            \item \textbf{Fleet Status}: État de santé de 20 agents en temps réel
            \item \textbf{Métriques}: Suivi de la précision et du recall en temps réel
        \end{itemize}
    \end{block}
    
    \vspace{0.5cm}
    
    \begin{center}
        {\Large\textbf{URL:} \texttt{http://localhost:8501}}
        
        \vspace{0.4cm}
        
        \begin{alertblock}{Demo Live}
            Lancement d'une simulation avec injection d'anomalies
        \end{alertblock}
    \end{center}
\end{frame}

\begin{frame}{Scénario de Démonstration}
    \textbf{Étape 1: Initialisation (T+0s)}
    \begin{itemize}
        \item Déploiement automatique de 20 agents
        \item Formation de la topologie mesh
        \item Entraînement initial des modèles River
    \end{itemize}
    
    \vspace{0.4cm}
    
    \textbf{Étape 2: Phase de Surveillance (T+10s)}
    \begin{itemize}
        \item Flux de données normal (non-anomalies)
        \item Apprentissage continu des patterns normaux
        \item Communication inter-agents active
    \end{itemize}
    
    \vspace{0.4cm}
    
    \textbf{Étape 3: Injection d'Anomalies (T+30s)}
    \begin{itemize}
        \item Spike de valeurs sur plusieurs capteurs
        \item Observation du processus de consensus
        \item Validation de la détection collective
    \end{itemize}
\end{frame}

\begin{frame}[fragile]{Processus de Consensus en Action}
    \begin{columns}[c]
        \begin{column}{0.55\textwidth}
            \begin{center}
                \textbf{Vote majoritaire}
                
                \vspace{0.3cm}
                
                \begin{tikzpicture}[scale=0.9]
                    % Agent central with anomaly detected
                    \node[circle, draw=red!70, line width=1.2pt, fill=red!15, minimum size=1.2cm, font=\normalsize\sffamily\bfseries] (center) at (0,0) {A5};
                    \node[below=0.1cm of center, font=\small, text=red!80] {Anomalie!};
                    
                    % Surrounding agents (only neighbors in ring)
                    \node[circle, draw=teal!70, line width=1.2pt, fill=teal!15, minimum size=1cm, font=\normalsize\sffamily\bfseries] (a4) at (-2,0) {A4};
                    \node[circle, draw=teal!70, line width=1.2pt, fill=teal!15, minimum size=1cm, font=\normalsize\sffamily\bfseries] (a6) at (2,0) {A6};
                    
                    % Ring connections
                    \draw[->, line width=1pt, color=teal!60] (a4) -- (center);
                    \draw[->, line width=1pt, color=teal!60] (center) -- (a6);
                    
                    % Vote labels
                    \node[above=0.05cm of a4, font=\small, text=green!70] {✓ OK};
                    \node[above=0.05cm of a6, font=\small, text=green!70] {✓ OK};
                \end{tikzpicture}
            \end{center}
        \end{column}
        
        \begin{column}{0.42\textwidth}
            \textbf{Résultat du Vote:}
            
            \vspace{0.3cm}
            
            \begin{block}{Décompte}
                \begin{itemize}
                    \item \textbf{2 voisins}: Pas d'anomalie
                    \item \textbf{1 agent (A5)}: Anomalie
                \end{itemize}
            \end{block}
            
            \vspace{0.3cm}
            
            \begin{alertblock}{Décision}
                \textbf{Faux Positif} détecté \\
                $\rightarrow$ Pas d'alerte
            \end{alertblock}
            
            \vspace{0.2cm}
            
            {\small Le consensus filtre le bruit local !}
        \end{column}
    \end{columns}
\end{frame}

% ===============================
% SLIDE 5: Résultats
% ===============================
\begin{frame}{Résultats Expérimentaux}
    \begin{table}
        \centering
        \begin{tabular}{lcc}
            \toprule
            \textbf{Métrique} & \textbf{Valeur} & \textbf{Objectif} \\
            \midrule
            \textbf{Recall} (Anomalies majeures) & \textbf{91.2\%} & > 90\% \\
            \textbf{Precision} & \textbf{62.4\%} & > 60\% \\
            F1-Score & 0.743 & > 0.70 \\
            \midrule
            Latence consensus moyenne & < 200ms & < 500ms \\
            Agents testés simultanément & 20 & 10+ \\
            \textbf{Réduction faux positifs} & \textbf{-37\%} & > -30\% \\
            \bottomrule
        \end{tabular}
    \end{table}
    
    \vspace{0.4cm}
    
    \begin{alertblock}{Résultat Clé}
        Le consensus distribué améliore drastiquement la précision du système !
    \end{alertblock}
\end{frame}

% ===============================
% SLIDE 6: Impact du Consensus
% ===============================
\begin{frame}{Impact du Consensus Distribué}
    \begin{columns}[T]
        \begin{column}{0.48\textwidth}
            \textbf{Sans Consensus:}
            \begin{itemize}
                \item Détection locale uniquement
                \item Faux positifs élevés
                \item Sensible au bruit
                \item Alertes excessives
            \end{itemize}
            
            \vspace{0.3cm}
            
            \begin{block}{Problème}
                Un capteur défectueux génère des alertes non fondées
            \end{block}
        \end{column}
        
        \begin{column}{0.48\textwidth}
            \textbf{Avec Consensus:}
            \begin{itemize}
                \item Validation par les pairs
                \item \textbf{-37\%} faux positifs
                \item Robustesse au bruit
                \item Haute confiance
            \end{itemize}
            
            \vspace{0.3cm}
            
            \begin{alertblock}{Solution}
                Le vote majoritaire filtre les anomalies isolées
            \end{alertblock}
        \end{column}
    \end{columns}
\end{frame}

% ===============================
% SLIDE 7: Applications
% ===============================
\begin{frame}{Applications Industrielles}
    \textbf{Domaines d'Application:}
    
    \vspace{0.3cm}
    
    \begin{itemize}
        \item \textbf{Smart Grids}: Surveillance de réseaux électriques
        \begin{itemize}
            \item Détection précoce de surcharges
            \item Milliers de transformateurs
        \end{itemize}
        
        \vspace{0.2cm}
        
        \item \textbf{IoT Industriel}: Maintenance prédictive
        \begin{itemize}
            \item Capteurs de vibration, température
            \item Prévention des arrêts de production
        \end{itemize}
        
        \vspace{0.2cm}
        
        \item \textbf{Smart Cities}: Gestion du trafic
        \begin{itemize}
            \item Détection d'embouteillages
            \item Monitoring qualité de l'air
        \end{itemize}
        
        \vspace{0.2cm}
        
        \item \textbf{Datacenters}: Détection d'intrusions réseau
    \end{itemize}
\end{frame}

% ===============================
% SLIDE 8: Conclusion
% ===============================
\begin{frame}{Conclusion et Perspectives}
    \textbf{Contributions:}
    \begin{itemize}
        \item Système multi-agents \textbf{fully distributed} pour détection d'anomalies
        \item Apprentissage en ligne \textbf{O(1)} avec River
        \item Consensus réduisant les faux positifs de \textbf{37\%}
        \item \textbf{Auto-réparation} validée expérimentalement
    \end{itemize}
    
    \vspace{0.4cm}
    
    \textbf{Perspectives:}
    \begin{itemize}
        \item Déploiement sur edge devices réels (Raspberry Pi)
        \item Algorithmes ML plus sophistiqués (Isolation Forest online)
        \item Topologie dynamique avec découverte automatique
        \item Validation avec partenaires industriels
    \end{itemize}
    
    \vspace{0.2cm}
    
    \begin{block}{Impact}
        Démonstration de la viabilité de l'\textbf{Edge AI collaborative} pour infrastructures critiques
    \end{block}
\end{frame}


% ===============================
% Thank You Slide
% ===============================
{
\setbeamertemplate{footline}{}
\begin{frame}[plain]
    \vfill
    \begin{center}
        {\color{accent}\rule{0.4\textwidth}{0.12cm}}\\[0.9cm]
        {\Huge\bfseries\color{primary} Merci pour votre attention}\\[0.6cm]
        {\Large\color{textcolor} Questions ?}\\[0.9cm]
        {\color{accent}\rule{0.4\textwidth}{0.12cm}}
    \end{center}
    
    \vspace{1.3cm}
    
    \begin{center}
        {\large\textbf{Abderrahman ESSEBYITY} \quad \textbf{\&} \quad \textbf{Hamza OUARIR}}\\[0.35cm]
        {\normalsize Encadré par: \textbf{Saif Islam Bouderba}}\\[0.4cm]
        {\normalsize Master SDA -- IA Distribuée}
    \end{center}
    
    \vfill
\end{frame}
}

\end{document}
